% !TeX root = ../paper.tex

\section{Conclusion and Outlook}
\label{sec:conclusion}

This thesis asked whether iterative static-analysis feedback can improve the
maintainability of LLM-generated Python without reducing functional correctness.
Across three loop configurations run on 63 MBPP+ tasks with a locally hosted
Qwen2.5-Coder-7B-Instruct model, the answer is a qualified yes. The loop yields a
statistically significant reduction in cognitive complexity, concentrated on the
tasks that have structural headroom, and it does so without a statistically
significant loss of correctness. The correctness-guarded configuration, which
treats the test suite as a continuous-integration quality gate, strengthens this
into a guarantee: it eliminates every correctness regression while retaining the
complexity benefit.

The study also delivers a cautionary methodological result. The Maintainability
Index does not improve and in fact declines slightly under every configuration, by
a similar amount whether or not comments are stripped. Because the baseline already
sits near the metric's ceiling and the Index conflates code size and comments with
structural quality~\citep{heitlager2007maintainability}, this decline is best read as evidence that the
Maintainability Index is an unsuitable outcome metric for this kind of
experiment rather than as a genuine loss of maintainability. Cognitive complexity
is the more informative signal, and it is the metric on which the qualified-yes
answer rests.

The results are bounded by a single model, a single seed, one decoding
configuration, and short benchmark tasks; the correctness test has limited
statistical power; and the reported magnitudes describe archived runs rather than
values a fresh execution will reproduce exactly. \Cref{sec:limitations} sets out
these bounds in full and derives the corresponding future work, of which the most
direct are replication across models, a volume-controlled maintainability measure,
human readability assessment, and a systematic study of how much the effect depends
on the wording of the feedback prompt.

What the study can offer in exchange for those bounds is traceability. Every
generated solution is retained with the measurements taken from it, and each figure
reported here can be recomputed from those artifacts, so the claims are open to
inspection at the level of the individual task rather than only in aggregate.
